\chapter{Инфляция, структура и утрата начальных условий}

\section{Омега и свидетельства в пользу плоскостности}

Вопрос студента о большой \(\Omega\) возвращает обсуждение к уравнению Фридмана—Робертсона—Уокера. Для современной эпохи, уже далеко после доминирования излучения, закон расширения записывают с членом вакуумной энергии и членом материи, тогда как излучение опускают, потому что сейчас оно очень мало:
\begin{equation}
H^2=\left(\frac{\dot a}{a}\right)^2=\frac{8\pi G}{3}\left(\rho_0+\rho_m\right)-\frac{k}{a^2}.
\end{equation}
Вакуумный член не разрежается. Член материи разрежается:
\begin{equation}
\rho_m=\frac{M}{a^3}.
\end{equation}
\(\Omega\) — это доля в уравнении расширения, приходящаяся на различные компоненты.

Разделим на \(H^2\). Тогда нормированные члены в сумме дают единицу. В приведенной к единообразию нотации
\begin{equation}
\begin{aligned}
\Omega_0&=\frac{8\pi G\rho_0}{3H^2},\\
\Omega_m&=\frac{8\pi G\rho_m}{3H^2},\\
\Omega_k&=-\frac{k}{a^2H^2},
\end{aligned}
\end{equation}
и, следовательно,
\begin{equation}
\Omega_0+\Omega_m+\Omega_k=1.
\end{equation}
\editorialnote{Звук и изображение доски на этом участке нечетки, и лектор переходит между несколькими названиями для вклада вакуума. Обозначения здесь приведены к единообразию для этого рассуждения после деления на \(H^2\).}
Если пространство плоское, \(k=0\), то
\begin{equation}
\Omega_0+\Omega_m=1.
\end{equation}
Если сумма материи и вакуума больше единицы, кривизна положительна. Если она меньше единицы, кривизна отрицательна.

Далее наблюдательная логика проста. Измерьте обычную материю и тёмную материю. Измерьте хаббловское расширение. Геометрически измерьте кривизну на очень больших масштабах. Заметьте также, что Вселенная ускоренно расширяется. Признаком вакуумного члена служит не \(H\) само по себе, а стремление \(\dot a/a\) к постоянной, а это значит, что масштабный фактор приближается к экспоненциальной ветви,
\begin{equation}
a(t)\propto e^{Ht}.
\end{equation}
Эти несколько измерений взаимно проверяют друг друга и показывают, что Вселенная очень близка к плоской, с точностью до нескольких процентов.

\begin{classroomqa}
\lecturetimestamp{00:10:50}
\audiencequestion{Включает ли материя тёмную материю?}
\lecturerresponse{Да. В рамках нынешнего обсуждения член материи включает тёмную материю. Можно было бы добавить и излучение, но сегодня эта доля настолько мала, что ею можно пренебречь.}
\end{classroomqa}

\begin{classroomqa}
\lecturetimestamp{00:11:11}
\audiencequestion{Есть ли какая-то мера тёмной энергии, кроме постоянной Хаббла?}
\lecturerresponse{Суть не в значении \(H\) самом по себе, а в том, как расширение стремится к экспоненциальному закону. Вместе с измерением материи и геометрическим измерением плоскостности это позволяет сделать вывод о вакуумном члене.}
\end{classroomqa}

\section{От плоскостности к структуре}

Почему Вселенная столь плоская? Предлагаемый ответ — инфляция. Плоская означает очень большая, очень гладкая и очень однородная на масштабах миллиардов, точнее десятков миллиардов, световых лет. Экспоненциальное расширение может растянуть область так, что она будет выглядеть именно так.

Но и хорошего может оказаться слишком много. Совершенно гладкая Вселенная остается совершенно гладкой. Гравитация не создает неоднородности из точной симметрии. Если структура существует сейчас, значит, раньше она должна была начаться с небольшой неоднородности. Если проследить нынешнюю Вселенную назад, то неоднородность ко времени прозрачности получается порядка одной части на \(10^5\):
\begin{equation}
\frac{\delta\rho}{\rho}\sim 10^{-5}.
\end{equation}
Вселенная была чрезвычайно гладкой, но не идеально.

Начиная с этого момента гравитация делает то, чего обычные силы в комнате, наполненной воздухом, не делают. Небольшая область повышенной плотности притягивает к себе материю. Соседняя область обедняется, так что область пониженной плотности становится ещё менее плотной, а область повышенной плотности — ещё более плотной. Гравитация усиливает контраст вместо того, чтобы стирать его.

Естественный вопрос состоит в том, почему этого не произошло раньше. Ответ — давление. Раньше давление излучения было велико. Используя
\begin{equation}
p=w\rho,
\end{equation}
получаем знакомые случаи
\begin{equation}
w_{\rm rad}=\frac13,\qquad
w_m=0,\qquad
w_{\rm vac}=-1.
\end{equation}
Положительное давление делает среду более однородной. Оно вытесняет возмущения и сглаживает среду. Материя, рассматриваемая на больших масштабах как не имеющая давления, такой опоры не даёт. Когда давление есть, гравитация сдерживается. Когда давлением можно пренебречь, побеждает гравитация.

Тёмная материя играет решающую роль в последовательности событий. Первыми схлопывающимися сгустками являются гало тёмной материи. Затем обычное барионное вещество падает в эти гало и, поскольку барионы могут диссипировать, собирается ближе к центру. Тёмная материя не диссипирует таким же образом, поэтому остаётся распределённой по всему гало. Начало этого сжатия происходит немного раньше эпохи прозрачности, примерно на порядок по времени, поэтому микроволновой фон не показывает его начало непосредственно.

Краткое отступление проясняет дальнейшую химическую историю. Почти все элементы тяжелее гелия — продукты звёздной эволюции, а не космологии ранней Вселенной. В ранней Вселенной возникают по существу водород, гелий и лишь малые следы нескольких лёгких ядер.

На небе наблюдается не сама плотность массы, а температурная анизотропия микроволнового фона. В неё вносят вклад флуктуации плотности. Вносят вклад и флуктуации скорости. Области повышенной плотности обычно выглядят немного более холодными, потому что фотоны теряют дополнительную энергию, выбираясь из более глубоких гравитационных ям, а доплеровские сдвиги из-за движения газа накладываются сверх этого. Полный численный анализ сложен. Надёжное утверждение проще: к тому времени, когда Вселенная стала прозрачной, плотность энергии менялась примерно на одну часть из \(10^5\).

\begin{classroomqa}
\lecturetimestamp{00:20:29}
\audiencequestion{Почему гравитация не заставила сгустки расти раньше?}
\lecturerresponse{Потому что раньше во Вселенной доминировало излучение, а излучение несло достаточно большое положительное давление, чтобы делать среду однородной. Только когда излучение стало несущественным по сравнению с материей, гравитация начала доминировать в росте этих сгустков.}
\end{classroomqa}

\begin{classroomqa}
\lecturetimestamp{00:24:35}
\audiencequestion{Разве для коллапса не нужен механизм диссипации?}
\lecturerresponse{Обычному барионному веществу действительно нужна диссипация, чтобы осесть к центру. Но тёмная материя может начать с образования гало и без такого рода диссипации при столкновениях. Есть и космологический вид затухания, связанный с самим расширением, и позже это окажется важным.}
\end{classroomqa}

\begin{classroomqa}
\lecturetimestamp{00:30:34}
\audiencequestion{Были ли также флуктуации скоростей, чтобы позднее галактики могли приобрести угловой момент?}
\lecturerresponse{Да. Первичная картина включала флуктуации плотности и флуктуации скоростей. На небе виден температурный рисунок, создаваемый и вариациями плотности, и движением газа.}
\end{classroomqa}

\section{Инфлатон и плато}

Теперь вводится механизм — скалярное поле \(\phi\). Его плотность энергии состоит из трёх частей:
\begin{equation}
\rho_\phi=\frac12\dot\phi^2+\frac12(\nabla\phi)^2+V(\phi).
\end{equation}
Есть вклад временного изменения. Есть вклад пространственного изменения. Скалярное поле может также нести настоящую потенциальную энергию \(V(\phi)\). Эта часть не экзотична; привычный пример — поле Хиггса. Здесь важна форма потенциала.

Потенциал должен быть довольно плоским на некотором интервале, с небольшим нисходящим наклоном, а затем при критическом значении дно резко уходит вниз. Он падает не точно до нуля. Он падает до чего-то очень малого. В единицах Планка современная вакуумная энергия имеет порядок
\begin{equation}
10^{-123}.
\end{equation}
Там, наверху на плато, она была гораздо больше этого, но всё же намного ниже планковского масштаба.

\begin{figure}[t]
\centering
\includegraphics[width=0.78\textwidth]{lecture_08_figure_03.png}
\caption{Качественный вид потенциала инфлатона с широким плато, крутым спадом и низким минимумом.}
\lectureframe{00:43:23}
\end{figure}

Этот график описывает не место в пространстве. Он описывает значение поля. Сказать, что Вселенная начиналась высоко на плато, значит лишь сказать, что значение поля в очень большой области лежало там, наверху. Неважно, если сначала \(\phi\) несколько менялось от места к месту, потому что инфляция чрезвычайно растягивает эти вариации. После достаточно большого расширения любой локальный участок выглядит почти однородным. \editorialnote{Подписи осей в кадре с доской неразборчивы, но устное обсуждение ясно показывает, что вертикальная ось — это потенциальная энергия \(V(\phi)\), а горизонтальная — значение поля \(\phi\).}

Пока поле находится на широкой, медленно спадающей части потенциала, потенциальная энергия ведёт себя как вакуумная энергия. Если бы \(\phi\) было полностью заморожено, оно вело бы себя в точности как вакуумная энергия. Если не учитывать опущенные численные множители, соотношение, сформулированное на лекции, просто имеет вид
\begin{equation}
H\propto \sqrt{V(\phi)},
\end{equation}
и это даёт почти экспоненциальный закон
\begin{equation}
a(t)\propto e^{Ht}.
\end{equation}
Почти всё, что теперь можно наблюдать, происходит после того, как поле переваливает через край. Спад — это место, где большая вакуумоподобная энергия на плато превращается в излучение, обычную материю, тёмную материю и крошечную остаточную вакуумную энергию, наблюдаемую сегодня.

\begin{classroomqa}
\lecturetimestamp{00:44:07}
\audiencequestion{На чём определено поле \(\phi\), и продолжают ли по-прежнему использоваться уравнения Эйнштейна?}
\lecturerresponse{Да. \(\phi\) — это поле на пространстве-времени, и то же уравнение Фридмана по-прежнему используется. В этом режиме плотность энергии в правой части задаётся главным образом \(V(\phi)\), если поле меняется достаточно медленно.}
\end{classroomqa}

\begin{classroomqa}
\lecturetimestamp{00:47:04}
\audiencequestion{Как вакуумоподобная энергия инфлатона может превратиться в материю и излучение?}
\lecturerresponse{Когда поле скатывается к минимуму и осциллирует, эти осцилляции и есть кванты поля. Если они нестабильны, то могут распадаться на фотоны, электроны, позитроны и другие знакомые частицы.}
\end{classroomqa}

\section{Медленное скатывание и хаббловское трение}

После того как поле переваливает через край, почему оно не должно просто колебаться вечно? Почему оно не должно качаться туда-сюда и, возможно, даже снова подняться вверх? Ответ, который здесь дается, состоит в космологическом механизме затухания. Это не настоящее микроскопическое трение, но в уравнении движения оно ведет себя именно так.

Пренебрежем пространственными градиентами и будем считать поле однородным. Тогда кинетический и потенциальный члены выглядят как одна степень свободы с динамикой \(T-V\), но, поскольку это плотности энергии, их нужно умножить на физический объем \(a^3\). Тем самым получаем однородный эффективный лагранжиан вида
\begin{equation}
L_{\mathrm{hom}}=a^3(t)\left[\frac12\dot\phi^2-V(\phi)\right].
\end{equation}
\editorialnote{Верхняя строка на доске частично обрезана. Этот однородный лагранжиан восстановлен лишь в той мере, в какой это необходимо, чтобы воспроизвести надежно читаемое уравнение движения, записанное под ним.}

Надежно читаемое на доске уравнение имеет вид
\begin{equation}
\frac{d}{dt}\left(a^3\dot\phi\right)=-a^3\frac{\partial V}{\partial\phi}.
\end{equation}
Раскрывая производную по времени, получаем
\begin{equation}
a^3\ddot\phi+3a^2\dot a\,\dot\phi=-a^3\frac{\partial V}{\partial\phi},
\end{equation}
а деление на \(a^3\) дает
\begin{equation}
\ddot\phi+3H\dot\phi=-\frac{\partial V}{\partial\phi}.
\end{equation}
Знак минус явно исправляется в устном выводе: сила равна минус производной потенциальной энергии.

\begin{figure}[t]
\centering
\includegraphics[width=0.78\textwidth]{lecture_08_figure_05.png}
\caption{Уравнение движения однородного поля инфлатона с членом хаббловского трения \(3H\dot\phi\).}
\lectureframe{00:57:10}
\end{figure}

Член \(3H\dot\phi\) и называется хаббловским трением. Он целиком возникает при дифференцировании множителя \(a^3\). Однако динамически он выглядит как член сопротивления в движении камня, падающего в воде. Здесь есть силовой член, член ускорения и член, линейный по скорости. Поле быстро достигает установившейся скорости и затем медленно движется вдоль плато. Это и есть режим медленного скатывания. Пока наклон пологий и поле движется медленно, потенциальная энергия меняется лишь медленно, так что расширение остается близким к экспоненциальному.

Вблизи минимума поле осциллирует. Эти осцилляции и есть кванты инфлатона. Если они нестабильны, то распадаются, и именно так энергия на плато превращается в горячее вещество и излучение более поздней Вселенной.

\begin{classroomqa}
\lecturetimestamp{00:50:48}
\audiencequestion{Почему поле просто не осциллирует вечно и что обеспечивает затухание?}
\lecturerresponse{Само расширение дает член, подобный трению. Множитель \(a^3\) в однородном лагранжиане порождает член \(3H\dot\phi\), и тот ведет себя как вязкое сопротивление.}
\end{classroomqa}

\section{Квантовые флуктуации и выход за горизонт}

Теперь возникает кажущаяся проблема. Если инфляция делает Вселенную чрезвычайно гладкой, почему она не стирает те самые зародыши, которые нужны галактикам? Ответ дает квантовая механика. Сохранившиеся флуктуации — не остатки несовершенного разглаживания. Они порождаются заново.

На мгновение забудем о расширении. Даже если бы поле было почти идеально однородным, квантовая механика не позволила бы точно зафиксировать одновременно \(\phi\) и \(\dot\phi\). Основное состояние все равно не лишено флуктуаций. По той же причине у обычного гармонического осциллятора есть нулевая энергия. В обычных условиях эти флуктуации трудно увидеть, потому что они коротковолновые и быстро осциллируют, так что наблюдения их усредняют.

Расширение меняет картину. В сопутствующих координатах мода поля сохраняет примерно ту же координатную длину волны, тогда как физическая длина волны растет вместе с масштабным фактором:
\begin{equation}
\lambda_{\rm phys}=a(t)\,\Delta x.
\end{equation}
Во время инфляции физический хаббловский масштаб приблизительно фиксирован,
\begin{equation}
d_H\sim H^{-1},
\end{equation}
поэтому соответствующий сопутствующий горизонт сжимается как
\begin{equation}
d_{H,\mathrm{coord}}\sim \frac{1}{aH}.
\end{equation}

Аналогия с веревкой делает это замораживание интуитивно понятным. Волна колеблется, потому что каждый участок веревки тянут соседние участки. Но если длина волны становится столь же большой, как горизонт, то существенные соседи оказываются вне причинного контакта. Восстанавливающее натяжение исчезает. Как только мода становится длиннее горизонта, она замораживается.

Это происходит снова и снова. Более длинные волны замораживаются раньше. Более короткие — позже. Даже мода, которая начинается с очень малой длины волны, в конце концов растягивается настолько, что пересекает горизонт. Таким образом, инфляция не просто стирает структуру. Она продолжает порождать замороженные вариации поля.

И этот процесс не прекращается просто потому, что инфляция длится долго. Старые замороженные моды растягиваются до огромных масштабов, но новые все время замораживаются. Накопленная картина становится статистически стационарной в смысле, существенном для космологии. В этом и состоит происхождение приблизительно масштабно-инвариантного спектра: один и тот же процесс пересечения горизонта повторяется на последовательных логарифмических масштабах и дает примерно одну и ту же амплитуду при переходе от одного диапазона длин волн к следующему.

После того как поле сходит с плато, горизонт снова растет. Моды, находившиеся вне временного инфляционного горизонта, позже возвращаются внутрь. В обычных терминах говорят, что во время инфляции они выходят за горизонт, а потом снова входят в него. К тому времени они растянуты до макроскопических масштабов, так что, когда они снова начинают осциллировать, это уже очень крупные структуры.

\begin{classroomqa}
\lecturetimestamp{01:12:02}
\audiencequestion{Как выглядит горизонт на простран\-ственно-временной диаграмме?}
\lecturerresponse{Физическое расстояние до горизонта во время экспоненциального расширения остается примерно постоянным и равно \(c/H\), но на сопутствующей диаграмме это фиксированное физическое расстояние выглядит все меньше и меньше, потому что координатный масштаб растягивается. В сопутствующих координатах горизонт сжимается как \(1/(aH)\).}
\end{classroomqa}

\begin{classroomqa}
\lecturetimestamp{01:15:59}
\audiencequestion{Что происходит с амплитудой по мере растяжения флуктуаций, и существует ли стационарное состояние?}
\lecturerresponse{Амплитуда отдельных мод уменьшается по мере их растяжения, но новые моды продолжают замораживаться. В результате возникает статистически стационарная, приблизительно масштабно-инвариантная картина, а не полное исчезновение флуктуаций.}
\end{classroomqa}

\section{От замороженных вариаций поля к вариациям плотности}

Как только поле приобретает небольшие вариации от места к месту, они превращаются в различия плотности, когда разные области переваливают через край в немного разные моменты времени. Область, находящаяся чуть дальше вниз по склону, первой достигает обрыва. Если ее вакуумная энергия превращается в излучение, это излучение разбавляется при расширении. Соседняя область, все еще находящаяся на плато, остается под доминированием вакуумной энергии и не разбавляется. К тому времени, когда падает вторая область, первая уже успевает потерять часть своей энергии из-за красного смещения.

Вот почему более медленное скатывание дает большие итоговые контрасты плотности. Если поле движется вдоль плато лишь очень медленно, то небольшая разница в значении поля соответствует большему запаздыванию по времени между тем, как одна область и следующая за ней достигают обрыва. У области, упавшей раньше, больше времени на то, чтобы разбавиться, прежде чем область, упавшая позже, преобразует свою энергию. Чем более полог склон, тем меньше установившаяся скорость и тем больше итоговая вариация плотности. Более высокое плато тоже обычно дает более крупные флуктуации. Ширина плато в основном определяет, как долго длится инфляция.

Данными о флуктуациях жестко ограничена лишь комбинация высоты и наклона. Если энергетический масштаб достаточно велик, должен быть также видимый вклад гравитационных волн. Численные границы приводятся лишь как грубые ориентировочные оценки. Грубо говоря, если бы плотность энергии на плато в планковских единицах была намного выше примерно \(10^{-14}\), микроволновый фон уже показывал бы слишком сильную структуру, обусловленную гравитационными волнами; если бы она была намного ниже примерно \(10^{-16}\), реалистических шансов увидеть такой сигнал было бы мало. \editorialnote{Говорящий явно рассматривает эти числа как приближенные и говорит, что не помнит точных значений. Они сохранены здесь только как утверждения о порядке величины.}

Единицы измерения объясняются обычным для натуральных единиц образом. Комптоновская длина волны частицы массы \(m\) равна
\begin{equation}
\lambda_C=\frac{\hbar}{mc}.
\end{equation}
Если \(c=\hbar=1\), то масса и обратная длина имеют одни и те же единицы измерения. Плотность энергии — это энергия в единице объема, поэтому ее размерность есть четвертая степень массы:
\begin{equation}
[\rho]=(\mathrm{mass})^4.
\end{equation}

Если смотреть на одну фиксированную точку, накопление замороженных мод выглядит как случайное блуждание. Разные точки пространства проходят разные случайные блуждания, и именно это превращает квантовый шум в пространственную вариацию.

\begin{classroomqa}
\lecturetimestamp{01:33:04}
\audiencequestion{Испытывает ли \(\phi\) случайное блуждание и делает ли это более раннюю начальную точку несущественной?}
\lecturerresponse{В любой фиксированной точке последовательно замораживающиеся моды складываются со случайной фазой, так что \(\phi\) действительно совершает случайное блуждание. Именно этот успех и составляет трудность: инфляция продолжает заменять более раннюю информацию более поздними флуктуациями, так что восстановить стадию до инфляции становится трудно.}
\end{classroomqa}

\section{Что скрывает инфляция}

Теория успешна в очень определённом смысле. Она связывает флуктуации в огромном диапазоне масштабов с одним и тем же лежащим в основе механизмом. Но тот же механизм сглаживает то, что было раньше. Более ранние неоднородности растягиваются и исчезают, а релевантными остаются квантовые флуктуации, возникшие позже. Поэтому инфляция одновременно и привлекательна, и разочаровывает. Она объясняет гладкость и спектр флуктуаций, но также стирает значительную часть более ранней истории.

Это вызывает тревогу по поводу отправной точки. Очень высокая плотность энергии — не проблема. Проблема в том, что энергия плато не выглядит ни естественно планковской, ни естественно крошечной. В обычных единицах она велика, но в планковских единицах всё же мала. Она не выглядит очевидной первой остановкой для фундаментальной теории.

В этой простой картине также нет буквального \(t=0\). Как только инфляция длится слишком долго, почти вся информация о любой более ранней стадии исчезает. Если бы число e-фолдингов было лишь немного выше необходимого минимума, возможно, сохранился бы какой-то след предыдущей стадии, но в общем случае инфляция стирает прошлое.

Та же осторожность появляется и при рассмотрении пространства де Ситтера. Чисто экспоненциальный закон верен для инфляционной ветви на поздних временах, но точное решение де Ситтера лучше описывается так:
\begin{equation}
a(t)\propto \cosh(Ht).
\end{equation}
Это решение сжимается до минимального размера, а затем расширяется. Следует ли понимать сжимающуюся ветвь в космологии буквально, остаётся открытым вопросом.

\begin{classroomqa}
\lecturetimestamp{01:37:31}
\audiencequestion{Могло ли не быть никакого начала, так что экспоненциальное поведение продолжалось бы неограниченно далеко в прошлое?}
\lecturerresponse{Не в этой простой картине. Точное пространство де Ситтера не является вечно расширяющейся экспонентой на всём протяжении времени; оно скорее имеет вид \(\cosh(Ht)\), со сжимающейся и расширяющейся ветвями. Следует ли понимать ту более раннюю ветвь буквально, оставляется открытым вопросом.}
\end{classroomqa}

\section{Подстройка, кривизна и антропные ограничения}

В первом приближении важными параметрами являются высота плато, его наклон и его ширина \(\Delta\phi\). Наклон важен потому, что он задаёт установившуюся скорость и, следовательно, величину флуктуаций. Ширина важна потому, что помогает определить, как долго длится инфляция. Распространённое предубеждение состоит в том, что характерный пробег поля не должен быть намного больше массы Планка, а плотность энергии не должна превышать величину порядка единицы в планковских единицах. При таких грубых ограничениях область параметров, дающая достаточную инфляцию, не занимает всё пространство параметров. Приведённая здесь оценка — подстройка на уровне нескольких процентов, около \(2\%\), чтобы сделать наклон достаточно пологим для примерно \(60\)–\(62\) e-фолдингов.

\begin{classroomqa}
\lecturetimestamp{01:46:02}
\audiencequestion{Можно ли понимать инфляцию в терминах пространства параметров и насколько сильна подстройка?}
\lecturerresponse{Да. В первом приближении рассматривают высоту, наклон и ширину плато. В лекции приводится грубая оценка, что для получения достаточного числа e-фолдингов требуется подстройка на уровне нескольких процентов, около \(2\%\).}
\end{classroomqa}

Эту умеренную подстройку затем сравнивают с гораздо более острой загадкой современной космологической постоянной, порядка \(10^{-123}\). Затем обсуждение переходит к антропным вопросам, потому что обе проблемы можно переформулировать, спросив, что требуется для формирования структур.

В ньютоновской картине положительная космологическая постоянная действует как универсальное отталкивание, пропорциональное расстоянию. Если бы она была намного больше, ей не пришлось бы разрывать галактики сегодня, чтобы оказаться губительной. Достаточно было бы, чтобы она была настолько велика, что первичные флуктуации \(10^{-5}\) вообще не смогли сконденсироваться в галактики. Приведённый пример состоит в том, что значение примерно на два порядка величины выше наблюдаемого уже дало бы такой эффект. В этом и состоит мысль, связанная с аргументом Вайнберга: космологическая постоянная не обязана быть в точности нулевой; достаточно, чтобы она была достаточно мала и допускала образование структур.

Отрицательная кривизна действует в том же направлении. На ньютоновском языке предыдущих лекций \(k<0\) означает разбегание быстрее скорости убегания. Слишком сильное устремление наружу тоже не позволяет малым первичным сверхплотностям когда-либо сконденсироваться в галактики. Инфляция разбавляет эту кривизну. Приведённая нижняя граница составляет около \(59.5\) e-фолдингов. Значение около \(62\) затем принимается как репрезентативная инфляционная история, заведомо лежащая выше порога.

Это меняет вопрос о вероятности. Если спрашивать о голом объёме пространства параметров, дающем достаточную инфляцию, ответ может составлять лишь несколько процентов. Если же обусловиться тем, что мы находимся во вселенной, где образуются галактики и существуют наблюдатели, способные задать этот вопрос, условная вероятность может стать гораздо больше. На этом обсуждение и заканчивается: не окончательным решением, а загадками тонкой настройки, их антропной переинтерпретацией и признанием того, что именно здесь и продолжаются большие споры.
